\section{模型检验与灵敏度分析}
\label{sec:validation}

\subsection{独立验算}

为避免"用同一套代码验证自己"，本文采用\textbf{信息隔离}的验算方案：另写独立程序，不复用求解代码中的任何函数，分别用\textbf{不同的算法}与\textbf{不同的模型表示}复核核心结果。

\textbf{（1）确定性核心的双算法交叉验算。} 对问题1，除主模型的 LP 外，另构造两种独立算法：（i）\textbf{母线侧变量布局的独立线性规划}，即用"母线侧充电量 $c_t$、母线侧放电量 $d_t$"替代电池侧变量，约束改写为 $g_t+d_t+P_t\Delta t=L_t\Delta t+c_t+s_t$、$E_t=E_{t-1}+\eta_cc_t-d_t/\eta_d$，变量的个数、界与约束矩阵均与主模型不同；（ii）\textbf{动态规划}，以 $20$ \si{kWh} 为步长把储电量 $[1200,10800]$ 离散为 $481$ 个状态，逐时段作 Bellman 递推。结果见 \cref{tab:verify}：独立 LP 的目标值 $35101.6$ 元与主模型完全一致（相对偏差 $7\times10^{-8}$），独立 DP 的目标值 $35198.4$ 元略高（$+0.276\%$），符合"DP 是原问题的受限离散解、其最优值必不低于 LP 最优值"的理论预期，验证了确定性核心的最优性与实现正确性。

\textbf{（2）守恒关系与约束核验。} 由 \cref{eq:identity} 检验全天购电量分解式，残差 $1.5\times10^{-11}$ \si{kWh}；由 $E_{144}=E_0$ 检验充放电量守恒 $\sum_t a_t=\sum_t b_t$，相对偏差 $2.0\times10^{-14}$；逐日重构储电量 $E_{144}=E_0+\sum_t a_t-\sum_t b_t$ 与结果文件记录值的最大偏差为 $1.1\times10^{-11}$ \si{kWh}。全部储电量满足 $[1200,10800]$ \si{kWh}，充放电量越界量不超过 $7\times10^{-13}$ \si{kWh}（数值误差量级）。

\textbf{（3）结果文件的独立复算。} 用独立实现的结算公式重算问题2--4 的紧急购电量与费用：问题2 的紧急购电量复算值 $275908$ \si{kWh} 与求解输出一致，计划购电费与紧急购电费的最大相对偏差为 $1.8\times10^{-14}$；问题4-2 在\textbf{实际电价}下的计划费、紧急费复算偏差为 $0$。同时直接读取 5 个结果文件核对其结构：行数（$334$ 天）、每行 $144$ 个时段值无缺失、"全天购电量 $=$ 行和"的最大相对偏差 $3.4\times10^{-16}$、"全天购电费 $=\sum_t p_tg_t$"的最大相对偏差 $5.1\times10^{-16}$、逐日 0:00 与 24:00 储电量偏差为 $0$。

\textbf{（4）信息上界关系的自洽性。} 完美信息模型的最优费用应不高于任何实际策略，仅完美光伏（或负载）信息的最优费用应不高于不调整策略，且完美信息下紧急购电费应为零。实测（单阶段边界口径、不计违约费，与 \cref{tab:p3-pol} 一致）：P5 $=12229643\le$ P3 $=21369475$；P5$\prime$ $=19249721\le$ P0 $=21252227$；A\_load $=16048332\le$ P0；完美信息下紧急购电费为 $4.7\times10^{-10}$ 元（数值零）。若改用含违约费的全账基，P5 为 $14324819$ 元，同样不高于任何实际策略。四条关系全部成立。

\begin{table}[H]
  \centering
  \caption{独立验算结果汇总（共 40 项检查，全部通过）}
  \label{tab:verify}
  \small
  \begin{tabular}{p{3.4cm}p{4.6cm}rrl}
    \toprule[1.5pt]
    检查类别 & 独立方法 & 本实现值 & 参照值 & 判定 \\
    \midrule[1pt]
    问题1 最优值 & 独立 LP（母线侧变量布局） & 35101.6 & 35101.6 & 相对偏差 $7\times10^{-8}$ \\
    问题1 最优值 & 独立 DP（$20$ \si{kWh} 网格） & 35198.4 & 35101.6 & $+0.276\%$（离散化间隙） \\
    全天购电量分解 & 独立重算 $\sum g_t$ 与恒等式 & 59482.70 & 59482.70 & 残差 $1.5\times10^{-11}$ \\
    充放电守恒 & 独立重算 $\sum a_t$、$\sum b_t$ & 18666.60 & 18666.60 & 相对偏差 $2.0\times10^{-14}$ \\
    问题2 紧急购电量 & 独立结算公式 & 275908 & 275908 & 偏差 $0$ \\
    问题2 计划/紧急费 & 独立结算公式 & $1.53\times10^7$ & $1.53\times10^7$ & 相对偏差 $1.8\times10^{-14}$ \\
    问题2 计划侧平衡 & 独立重算带裕度预测后逐时段核验 & $\le10^{-4}$ & $0$ & 最大残差 $<10^{-4}$ \si{kWh} \\
    问题2 供电保障 & 逐时段检验 $g+\eta_db+P\Delta t+e\ge L\Delta t$ & $0$ & $0$ & 违反次数 $0$ \\
    结果文件结构 & 独立读取 5 个结果文件 & $334$ 行 & $334$ 行 & 行和与合计偏差 $\le5\times10^{-16}$ \\
    逐日储电量自洽 & 独立重构 $E_{144}$ & --- & --- & 最大偏差 $1.1\times10^{-11}$ \si{kWh} \\
    问题4 结算复算 & 按实际电价独立复算 & $1.38\times10^7$ & $1.38\times10^7$ & 相对偏差 $0$ \\
    信息上界关系 & 四条不等式与零紧急费检验 & --- & --- & 全部成立 \\
    \bottomrule[1.5pt]
  \end{tabular}
\end{table}

\subsection{灵敏度分析}

除问题1 的参数灵敏度（\cref{tab:p1-sens}）外，本文另作两项针对机制参数的检验。

\textbf{（1）问题2：紧急购电费率倍数的敏感性。} 保持计划不变、仅改变紧急购电费率倍数 $K$（题目给定 $K=5$）重新结算，结果见 \cref{tab:sens-k}。总费用对 $K$ 近似线性：$K$ 由 $5$ 降为 $3$ 时总费用下降 $3.2\%$，升为 $8$ 时上升 $4.7\%$；紧急购电费占比由 $5\%$ 升至 $12\%$。该结果说明\textbf{安全裕度的价值与紧急费率正相关}：若紧急费率更低，则最优策略会减少安全储备、更多依赖紧急购电。

\begin{table}[H]
  \centering
  \caption{问题2 对紧急购电费率倍数的敏感性（334 天合计）}
  \label{tab:sens-k}
  \small
  \begin{tabular}{ccccc}
    \toprule[1.5pt]
    紧急费率倍数 $K$ & 计划购电费 (元) & 紧急购电费 (元) & 总费用 (元) & 紧急费占比 \\
    \midrule[1pt]
    3 & 15347790 & 780117  & 16127907 & $4.8\%$ \\
    4 & 15347790 & 1040156 & 16387946 & $6.3\%$ \\
    \textbf{5} & 15347790 & 1300195 & \textbf{16647985} & $7.8\%$ \\
    6 & 15347790 & 1560234 & 16908024 & $9.2\%$ \\
    8 & 15347790 & 2080312 & 17428102 & $11.9\%$ \\
    \bottomrule[1.5pt]
  \end{tabular}
\end{table}

\textbf{（2）问题3：违约费率口径的敏感性。} 题目给出的费率结构等价于对称罚金 $0.5p_t$（\cref{eq:p3-sym}）。若改为更强的对称罚金 $1.0p_t$（即调低与调高分别按 $1.0$ 与 $2.0$ 倍计价），并重新求解各阶段的调整计划，结果见 \cref{tab:sens-dev}。

\begin{table}[H]
  \centering
  \caption{问题3 对违约费率口径的敏感性（全年总费用，单位：元）}
  \label{tab:sens-dev}
  \small
  \begin{tabular}{lcccc}
    \toprule[1.5pt]
    违约费率口径 & P0（不调整） & P1（$+$6:00） & P2（$+$12:00） & P3（题目方案） \\
    \midrule[1pt]
    基准（调低 $0.5$、调高 $1.5$） & 21252227 & 21421999 & 21369474 & 21369475 \\
    对称（调低 $1.0$、调高 $2.0$） & 21252227 & 21517036 & 21531996 & 21531999 \\
    \bottomrule[1.5pt]
  \end{tabular}
\end{table}

\begin{table}[H]
  \centering
  \caption{预报误差放大情景下的策略对比（334 天，单位：元）}
  \label{tab:sens-fc}
  \footnotesize
  \begin{tabular}{ccccccc}
    \toprule[1.5pt]
    预报误差倍数 & P0 & P1 & P2 & P3 & P3 相对 P0 & 12:00 更新边际价值 \\
    \midrule[1pt]
    1.0（原始） & 21252227 & 21421999 & 21369474 & 21369475 & $+0.55\%$ & 52525 \\
    1.5（放大） & 22294529 & 31526865 & 28488706 & 24518005 & $+9.97\%$ & 3038159 \\
    \bottomrule[1.5pt]
  \end{tabular}
  \par\vspace{3pt}{\footnotesize 注：放大方式为 $\tilde P=F^{\text{预报}}+1.5\,(F^{\text{预报}}-F^{\text{实际}})$，即把附件3 预报误差整体放大 1.5 倍后重做全部策略。}
\end{table}

两种口径下\textbf{结论一致且方向相同}：调整策略的费用均高于不调整的 P0，且费率越重、调整的劣势越明显（P3 相对 P0 的差额由 $+0.55\%$ 扩大到 $+1.32\%$）。这从机制层面确认了 \S5.6 的结论不依赖于某一特定费率取值。

\subsection{稳健性与口径讨论}

\textbf{（1）终端电量口径。} 本文取日周期条件 $E_{144}=E_0=6000$ \si{kWh}。问题2、3 中电价每天相同，稳态下跨日搬运无额外收益，该口径与跨日最优解的差异为二阶量级；问题4 中电价逐日波动，日周期会在一定程度上低估"低价日多存、高价日多用"的跨日套利收益，故问题4 的费用结果偏保守。作为对比，若允许终端电量自由，单日优化会把电量放空至下限，形成逐日退化，反而更不合理，故本文不采用。

\textbf{（2）违约费基准口径。} 本文以 0:00 计划为唯一基准一次结算。若改为"每次调整相对上一次承诺"，由三角不等式 $|\tilde g_t-\hat g_t|\le\sum_k|\Delta_t^{(k)}|$，链式口径的违约费不低于单基准口径；而本题违约费（$30.4$ 万元）相对紧急购电费（$742.8$ 万元）很小，即使翻倍也不改变"预报更新无正收益"的结论。

\textbf{（3）预报对齐口径。} 附件3 的"预报 $k$ 小时"必须按"发布时刻之后的第 $k$ 小时"理解。若误按"当日第 $k$ 小时"对齐，将产生 $6$--$18$ 小时的系统性错位（例如把 6:00 发布、面向 13:00--24:00 的预报误用于 7:00--18:00），据此得到的计划会严重偏离实际，结论将完全反转。本文在 \cref{tab:fc-acc} 中以"同一时间窗比较"的方式验证了本文口径的正确性：该口径下后续发布时刻的预报确实更准确，符合预报时效的物理规律。

\textbf{（4）主要误差来源。} 由信息价值分解（\cref{fig:p3-rev}(b)），负载信息价值（$520.39$ 万元）为光伏信息价值（$200.25$ 万元）的 $2.6$ 倍。因此本文模型的精度上限主要由负载预测决定；若引入负荷预报产品（气象驱动的短期负荷预测），问题2--4 的费用仍有较大改善空间。
